\exercisesfor{8}

\begin{enumerate}
\def\labelenumi{\arabic{enumi}.}
\item
  从  到  的映射 \(z \mapsto \bar{z}\) 是复李群 \(\mathbf{C}\) 的非解析连续自同构.\(\mathbf{C}\)\(\mathbf{C}\)
\item
  令 G 为由实数序列 \((\lambda_{1}, \lambda_{2}, \ldots)\) 组成的 Hilbert 空间, 满足 \(\sum_{i} \lambda_{i}^{2} < +\infty\). 将 G 视为实李群. 令 \(G_{n}\) 为由 \((\lambda_{1}, \lambda_{2}, \ldots) \in G\) 组成的集合, 其中对  有 \(\lambda_{m} \in \frac{1}{m} Z\). 各 \(1 \leqslant m \leqslant n\)\(G_{n}\) 是 G 的闭李子群, 因而 \(H = \bigcap_{n} G_{n}\) 是 G 的闭子群. 但此群全不连通且非离散, 因而不是 G 的李子群.
\end{enumerate}

¶ 3. 在 \(Q_{p} \times Q_{p}\) 中, 每个闭集都可以由一族解析方程定义. 推出定理 2 的推论 2 (ii) 在省去 I 有限的假设时不再成立.

¶ 4. 设 G 是有限维实李群, A 是 G 的一个子群. 称 L(G) 中的元素 x 是 A-可达的, 如果对 e 的每个邻域 U, 都存在从 \(\alpha\){[}0, 1{]} 到 A 的连续映射 , 使 \(\alpha(0) = e\), 且对  有 \(\alpha(t) \in \exp(tx)\) . U. 令 b 为 L(G) 中 A-可达元素的集合.\(0 \leqslant t \leqslant 1\)

\begin{enumerate}
\def\labelenumi{(\alph{enumi})}
\item
  证明 \(\mathfrak{h}\) 是 \(\mathbf{L}(\mathbf{G})\) 的李子代数. (使用 §6 的练习 19.)
\item
  令 H 是 G 的积分子群, 满足 \(\mathrm{L}(\mathrm{H}) = \mathfrak{h}\) . 证明 \(H \subset A\) . (令 I = \(\{-1, 1\}\), 令 \(R'\) 赋予欧氏范数. 令 \((x_{1}, \ldots, x_{r})\) 为 h 的一个基. 构造从 I 到 A 的连续映射 \(\alpha_{1}, \ldots, \alpha_{r}\) 以及从  到 R 的连续映射 \(f_{1}, \ldots, f_{r}\), 使得对 \(I'\) 中的 \(t = (t_{1}, \ldots, t_{r}) \in I'\) 有
\end{enumerate}

\[
\begin{array}{c} \alpha_ {1} (t _ {1}) \ldots \alpha_ {r} (t _ {r}) = \exp (f _ {1} (t) x _ {1}) \ldots \exp (f _ {r} (t) x _ {r}) \\ \| t - (f _ {1} (t), \ldots , f _ {r} (t)) \| \leqslant \frac {1}{2}. \end{array}
\]

然后应用以下定理: 令 \(f\) 为从 \(\mathbf{I}^r\) 到 \(\mathbf{R}^r\) 的连续映射, 使对所有  有 \(\| f(x) - x \| \leqslant \frac{1}{2}\); 则 \(x \in \mathbf{I}^r\)\(f(\mathbf{I}^r)\) 包含 \(\mathbf{R}^r. \dagger\) 中 0 的一个邻域

\begin{enumerate}
\def\labelenumi{(\alph{enumi})}
\setcounter{enumi}{2}
\item
  推出 \(\mathfrak{h}\) 是 A 在 e 处相切的子代数.\(e\)
\item
  证明若 A 是弧连通的, 则 A = H.†
\end{enumerate}

¶ 5. 设 G 是 Hausdorff 拓扑群, H 是 G 的闭子群, \(\pi\) 是从 G 到 G/H 的典范映射. 假设 H 是有限维实李群. 则 G/H 中 \(\pi(e)\) 的某个邻域 U 上存在连续映射 \(\sigma\) 到 G, 使 \(\pi \circ \sigma = \mathrm{Id}_{\mathrm{U}}\). (令 \(\rho\) 为 H 在 \(\mathbf{GL}(n,\mathbf{R})\) 上的解析线性表示, 它是局部同胚 (§ 6, 定理 1 的推论). 令 f 为 G 上满足 \(f\)\(\geqslant 0\) 的连续函数, 在 e 处等于 1, 并在 e 的某个充分小邻域外为零. 令 \(e\)\(e\)\(ds\) 为 H 的左 Haar 测度. 对 \(x \in G\), 记

\[
g (x) = \int f (x s) \rho (s) ^ {- 1} d s \in \mathbf {M} _ {n} (\mathbf {R}).
\]

于是对 \(g(xt) = g(x)\varphi(t)\)\(x \in G\) 和 \(t \in H\) 有 . 若 V 足够小,

\[
g (x) \in \mathbf {G L} (n, \mathbf {R})
\]

当 x 充分接近 e 时, \(x\). 最后使用如下事实: 要证明的定理对 \(e\)\(\mathbf{GL}(n, \mathbf{R})\) 和 \(\varphi(\mathbf{H})\) 在局部成立.

\begin{enumerate}
\def\labelenumi{\arabic{enumi}.}
\setcounter{enumi}{5}
\item
  设 G 是一个 C＾r 类群 (§ 5 练习 1), 其中 \(C^{r}\)\(r \geqslant 2\). 在 G 上存在唯一的实李群结构 S, 使 S 的 C＾r 类底层流形结构就是给定结构. (S 的唯一性由定理 1 的推论 1 得出. 令 L(G) 为 § 5 练习 2 中与 G 相关的可赋范李代数. 存在一个实李群芽 \(C^{r}\)\(L(G)\)\(G^{r}\) 以及从 \(L(G^{\prime})\) 到 L(G) 的同构 h (§ 4 定理 3). 仿照 § 4 第 1 号验证: 存在 \(L(G)\) 中  的对称开邻域 \(G^{\prime\prime}\), 以及从 \(e_{G}\) 到 G 的 C＾r 类映射 \(G^{r}\)\(\phi\), 满足 \(C^{r}\), 且对 \(G^{\prime\prime}\) 中的 \(T_{e}(\phi) = h\) 有 \(\phi(g_{1}g_{2}) = \phi(g_{1})\phi(g_{2})\). 缩小 \(g_{1}, g_{2}\)\(G^{\prime\prime}\)\(G^{\prime\prime}\) 后, 可假设 \(V = \phi(G^{r})\) 在 G 中开, 且 \(\phi\) 是从流形 \(C^{r}\)\(G^{\prime\prime}\) 到流形 V 的 C＾r 类同构. 因而 V 上存在一个实李群芽结构, 使其 C＾r 类底层流形结构为给定结构. 对所有 \(C^{r}\)\(g \in G\), Int g 定义了从 \(V \cap (g^{-1}Vg)\) 到 \((gVg^{-1}) \cap V\) 的解析映射 (定理 1). 根据 § 1 的命题 18, G 上存在实李群结构 S, 在 e 的某个开邻域上与 V 诱导相同的解析结构. 通过平移, S 在 G 上的 C＾r 类底层流形结构就是给定结构.)\(C^{r}\)
\item
  设 \(G\) 是实李群, H 是 G 的闭子群.\(H\)\(G\)
\end{enumerate}

\begin{enumerate}
\def\labelenumi{(\alph{enumi})}
\tightlist
\item
  令 \(\mathfrak{h}\) 为所有满足对  都有  的 \(x\in \dot{\mathbf{L}} (\mathbf{G})\) 所成的集合. 则 \(\exp (tx)\in \dot{\mathbf{H}}\)\(t\in \mathbf{R}\)\(\mathfrak{h}\) 是 \(\mathrm{L}(\mathrm{G})\) 的李子代数. (使用 §6 的命题 8.)
\end{enumerate}

\begin{enumerate}
\def\labelenumi{(\alph{enumi})}
\setcounter{enumi}{1}
\tightlist
\item
  假设 \(\mathbf{H}\) 是局部紧的. 证明 \(\mathbf{H}\) 是 \(\mathbf{G}\) 的李子群. (首先证明 \(\mathfrak{h}\) 是有限维的, 方法是证明 \(\mathfrak{h}\) 中存在一个相对紧邻域. 然后仿照定理 2 的证明, 选取 \(\mathrm{V}_2\), 使 \((\exp \mathrm{V}_2) \cap \mathrm{H}\) 相对紧.)
\end{enumerate}

